\hypertarget{chapter-2-why-latex-works-the-way-it-does}{%
\chapter{Why LaTeX Works the Way It Does}\label{chapter-2-why-latex-works-the-way-it-does}}

TeX and LaTeX are often treated as synonyms, but the distinction between them is the key to understanding why the system behaves the way it does, including the parts that feel inconsistent until the distinction is made explicit. TeX is Donald Knuth's typesetting engine~\cite{knuth1984texbook}: a fixed, extremely stable set of primitives for boxes, glue, penalties, fonts, and macro expansion, essentially unchanged since the 1980s. LaTeX is a large body of macros, written in TeX's own macro language, that sits on top of those primitives and gives them names and conventions a document author can use without engaging the primitives directly. \texttt{\textbackslash{}section}, \texttt{\textbackslash{}begin\{itemize\}}, and \texttt{\textbackslash{}cite} are not part of TeX; they are LaTeX macros that eventually expand down into TeX primitives such as box construction, penalty insertion, and font selection.

\hypertarget{content-and-convention}{%
\section{Content and Convention}\label{content-and-convention}}

This layering was a deliberate design choice, not an accident of history. Knuth built TeX to be a stable, low-level typesetting engine and explicitly left higher-level document conventions to be built separately, on top of it, by whoever needed them. Leslie Lamport's LaTeX~\cite{lamport1994latex} was one such layer, built to let an author write in terms of document structure (sections, theorems, references, bibliographies) rather than in terms of layout primitives (boxes, glue, explicit spacing). The practical consequence is that a LaTeX command is almost never a single behavior; it is a name for a bundle of lower-level TeX operations that LaTeX's author decided should happen together.

This is why two different LaTeX installations, or two different versions of a class file, can render the same source differently even though neither one is ``wrong'': the primitives TeX exposes are stable, but the conventions LaTeX builds on top of them are a design decision, open to revision, and frequently revised. A superuser benefits from being able to look past the convention to the primitive underneath it when a convention does not do what is needed, rather than searching for a package that happens to already provide the exact combination of primitives required.

\hypertarget{macro-expansion-is-rewriting-not-procedure-calling}{%
\section{Macro Expansion Is Rewriting, Not Procedure Calling}\label{macro-expansion-is-rewriting-not-procedure-calling}}

It is tempting, especially for anyone coming from a conventional programming background, to read a LaTeX document as a sequence of procedure calls: \texttt{\textbackslash{}section\{Introduction\}} looks like a function invocation, and it is natural to imagine TeX ``running'' the document from top to bottom the way an interpreter runs a script. The actual model, introduced in the previous chapter and worth restating here in this context, is closer to term rewriting. A macro is a rule that says: wherever this pattern of tokens appears, replace it with that sequence of tokens, and keep rewriting until no more expandable tokens remain in the position currently being processed.

The practical effect of this distinction shows up constantly in intermediate LaTeX use. A macro cannot, in general, ``return early'' the way a function can, because there is no call stack to return from, only a token stream being progressively rewritten. Recursion in TeX macros is not a language feature bolted on top of an imperative core; it falls directly out of the rewriting model, since a macro is permitted to expand to a token sequence that includes its own name, and expansion simply continues until some conditional stops it. Conditionals themselves (\texttt{\textbackslash{}ifnum}, \texttt{\textbackslash{}ifx}, and the LaTeX3 boolean and predicate machinery built on top of them) are not control-flow statements in the procedural sense; they are themselves expandable, producing one branch's tokens or the other's, which then continue to be rewritten like anything else.

Treating macros as rewrite rules rather than procedures also explains a class of bugs that otherwise seem inexplicable: a macro whose definition looks correct but produces the wrong output because it was invoked in a context where full expansion had already happened, or had not yet happened, by the time TeX needed a fully expanded value. Diagnosing this kind of failure is a matter of tracing where in the rewriting process a given token sequence currently sits, not a matter of stepping through a call stack.

\hypertarget{where-latex-breaks-the-abstraction}{%
\section{Where LaTeX Breaks the Abstraction}\label{where-latex-breaks-the-abstraction}}

No layer built on top of a lower one hides that lower layer perfectly, and LaTeX is no exception. There are places where the document-structure abstraction LaTeX offers leaks, and a superuser needs to recognize these leaks rather than be surprised by them.

Fragile commands are the clearest example. Certain LaTeX constructs, particularly those involving arguments that get written to an auxiliary file or reprocessed in a different context (section titles going into the table of contents, captions going into a list of figures), are not fully expanded in the context where they are written; they are expanded again later, in a different context, when the auxiliary file is read back in. A command that behaves correctly the first time but breaks when it is reprocessed this way is described as fragile, and \texttt{\textbackslash{}protect} exists specifically to suppress expansion at the first pass so the correct behavior happens at the second. This is not a design flaw so much as a visible seam between the token-rewriting substrate and the document-structure convention layer built on top of it.

Package interaction is a second seam. Because LaTeX macros can redefine other macros, and because the order in which packages are loaded determines which definition wins when two packages try to modify the same behavior, the outcome of loading a set of packages is not simply the union of what each package does in isolation; it depends on load order and on which package's redefinition happens to be the one still in effect when the document is typeset. This is a direct consequence of building convention out of the same rewriting mechanism an author's own macros use: packages are not sandboxed from each other or from the document.

A third seam is engine dependence. Primitives available in one engine (\texttt{\textbackslash{}pdfximage} in classic pdfTeX, direct OpenType font loading in LuaLaTeX and XeLaTeX) are not available in others, so a document that reaches past LaTeX's abstraction layer to use an engine-specific primitive stops being engine-agnostic at exactly that point, even though the surrounding LaTeX conventions remain the same.

None of these seams are reasons to avoid working close to the primitive layer. They are, if anything, the reason to understand it: a superuser who knows where LaTeX's abstraction is thin can predict which constructs will misbehave under reprocessing, package interaction, or engine change, and can either avoid them or work around them deliberately, rather than discovering the seam by accident in a five-hundred-page document three days before a deadline.
